% !TeX root = ../深度学习与大语言模型技术原理.tex
% 由主文档通过 \input 载入；请编译主文档。
\setcounter{chapter}{6} % 保留原书章号。
\chapter{数据治理、配比与合成数据基础}
\label{ch:data}

\section{数据处理是一串改变分布的操作}

预训练数据并不是抓取文本后简单拼接。解析、语言识别、文档过滤、去重、敏感字段处理和采样都会改变训练分布。一个可追溯数据单元应保存来源、时间、许可信息、处理版本、父文档标识和保留原因；对敏感原始字段则应采取符合任务约束的访问与存储方式。保留处理元数据，才能知道模型问题来自原始语料、解析错误还是过滤器偏差。

语言识别和质量分类器输出通常是估计分数，而非无误标签。阈值过严可能误删方言、代码混合文本和低资源语言；阈值过松又会保留乱码、模板垃圾和不适宜内容。文档级过滤适于处理整体来源和结构，词元或片段级过滤适于标记局部内容；直接删除局部词元也可能破坏语法与上下文，需要根据目标选择遮盖、剔除整段或保留标记。

个人身份信息（personally identifiable information, PII）处理可结合规则、实体识别与人工抽检。规则擅长固定格式，模型擅长依上下文识别，但都会误报或漏报。数据质量报告应给出各阶段输入量、保留率、抽检误差与语言领域分布，避免用一个“高质量分数”掩盖覆盖损失。

\section{精确去重、MinHash 与局部敏感哈希}

精确去重对规范化后的文档计算哈希，再核对相同键；过强规范化可能把本应不同的程序或数字合并。近重复检测常把文本变为 $n$-gram 集合 $A$，用 Jaccard 相似度 $J(A,B)=|A\cap B|/|A\cup B|$ 衡量重叠。重复转载、模板与局部抄录的判断还需结合文档长度和片段位置。

MinHash 为集合元素选一个随机排列，记录排列下最小元素。对非空集合 $A,B$，两个最小值相同，当且仅当 $A\cup B$ 中排名最小的元素属于交集。每个并集元素成为最小者的概率相同，故
\[
\Pr\{h_{\min}(A)=h_{\min}(B)\}=J(A,B).
\]
多个独立签名的碰撞比例可估计相似度。实际哈希近似随机排列，应以验证数据检查其误差。

\emph{局部敏感哈希（locality-sensitive hashing, LSH）}通过签名分桶缩小候选比较范围。若签名按 $b$ 个带、每带 $r$ 行分组，在理想独立假设下，相似度为 $s$ 的文档至少一个带完全相同的概率为 $1-(1-s^r)^b$。这说明带数与行数控制召回和误候选，而不是实现一个绝对相似度阈值。候选生成后还要进行更精确核对，且最好按重复簇而非单文档切分训练和测试集。

\section{领域混合、课程与污染检测}

设第 $j$ 个领域分布为 $\mathcal D_j$，按概率 $\alpha_j$ 抽样，则目标为
\[
\mathcal L(\theta;\alpha)=\sum_{j=1}^J\alpha_j\mathbb E_{Z\sim\mathcal D_j}[\ell(\theta;Z)],
\qquad \alpha_j\ge0,\quad \sum_j\alpha_j=1.
\]
这里的概率必须说明是按文档还是按词元定义。长文档较多的领域，即使文档抽样比例相同，词元贡献也可能更大。重要性重加权能够改变目标，但过大的权重可能提高梯度方差。

数据配比没有独立于模型和目标的通用最优比例。可以在代理模型上比较配比，进行等词元预算消融，再用较大模型验证是否迁移。课程学习（curriculum learning）让分布随训练进度变化，例如从短上下文扩展到长上下文；难度评分本身可能有偏，且“先易后难”不必然优于随机混合。

污染检测（contamination detection）检查训练数据与评测题目、答案、解析或近似改写的重叠。时间切分、来源隔离、字符串匹配与语义近邻各有盲点。应在模型调参前冻结测试数据，避免用测试反馈反复设计过滤器；否则即使训练语料没有原题，也可能通过实验选择泄漏测试信息。

\section{指令合成与迭代自训练}

合成数据（synthetic data）用模型生成任务、回答或反馈。Self-Instruct 风格流程从种子任务扩展指令，再生成响应，最后筛选格式、重复和质量；教师--学生生成则以较强教师提供目标。教师必须获得完成任务所需的真实证据，否则流畅回答仍可能传播错误。

筛选改变的不只是正确率，还包括题目难度、回答长度和风格。若只保留模型已能完成的题目，难题可能被系统性丢弃；若用同源模型同时生成与裁判，相关错误可能被反复强化。迭代自训练应保留独立人工或真实数据锚点，并报告每轮生成量、通过率、多样性和独立测试表现。其概率视角、蒸馏与拒绝筛选在第~\ref{ch:distillation} 章进一步展开。

